\block{Resumen}{
	Este estudio analiza la dinámica de producción en los 78 municipios de Puerto Rico durante el periodo (2021-2024).
	Utilizando un enfoque de econometría Bayesian, se una función de producción Cobb-Douglas jerárquica que integra el
	consumo energético como indicador de actividad económica. La metodología emplea Procesos Gaussianos en el Espacio de Hilbert
	(HSGP) para aislar choques macroeconómicos globales de las tendencias de crecimiento locales. Los resultados revelan una
	elasticidad de producción balanceada ($\approx 0.3$ para capital y $\approx 0.7$ para trabajo) e identifican disparidades
	significativas en la Productividad Total de los Factores (TFP) entre municipios. Esto revela que aunque la producción de
	Puerto Rico sigue la teoría la estimación de la producción de municipios es imprecisa.
}




